\chapter{GİRİŞ}
\label{chap:giris}

CRISPR-Cas9 sistemi, programlanabilir RNA rehberliğiyle DNA üzerinde belirli bölgelerin hedeflenebilmesini sağlayan temel bir genom düzenleme yaklaşımıdır. Cas9 enziminin RNA kılavuzlu hedef tanıma mekanizması, genom düzenleme çalışmalarında hedef bölge seçimini hesaplanabilir bir probleme dönüştürmüştür (Jinek ve diğ., 2012; Cong ve diğ., 2013). Bununla birlikte hedef sekansa benzeyen farklı genomik bölgelerde istenmeyen kesimlerin oluşabilmesi, CRISPR-Cas9 uygulamalarında hedef dışı bölge (off-target site) tahminini merkezi bir güvenilirlik sorunu haline getirmiştir (Hsu ve diğ., 2013).

Rehber RNA ve aday hedef DNA dizisi arasındaki eşleşme kalitesi, hedef dışı tahmin probleminin temel girdilerinden biridir. Aynı veya benzer sekans örüntüleri farklı kromatin bağlamlarında farklı erişilebilirlik ve aktivite davranışı gösterebilir. Cas9 bağlanması ve kesim etkinliğinin kromatin erişilebilirliği, nükleozom konumu ve hücresel bağlamla ilişkili olabileceği deneysel ve derleme çalışmalarında gösterilmiştir (Wu ve diğ., 2014; Horlbeck ve diğ., 2016; Isaac ve diğ., 2016; Yarrington ve diğ., 2018; Verkuijl ve Rots, 2019). Bu nedenle aday hedef bölgenin epigenetik durumu, nükleozom ilişkili özellikleri, deneysel ölçüm koşulu ve hedef gözleminin bağlamı tahmin probleminde ayrı bir bilgi ekseni oluşturur. crisprSQL veri tabanı CRISPR/Cas hedef dışı kesim deneylerini düzenli bir kaynak altında toplamış, Mak ve diğ. (2022) ise bu veri çizgisi üzerinde deneysel epigenetik ve hesaplanmış nükleozom tanımlayıcılarının hedef dışı aktiviteyle ilişkisini incelemiştir (Störtz ve Minary, 2021; Mak ve diğ., 2022).

Bu tezde incelenen problem, CRISPR-Cas9 hedef dışı tahmininde epigenetik ve kromatin bağlamının çizge tabanlı modeller içinde nasıl temsil edilmesi gerektiğidir. Çizge sinir ağları (graph neural networks, GNN), düğümler ve kenarlar arasında bilgi aktarımı yaparak ilişkisel yapıların öğrenilmesine olanak verir. Kipf ve Welling (2017) tarafından yaygınlaştırılan graph convolutional network (GCN) yaklaşımı, bu tür ilişkisel temsiller için temel bir model ailesi oluşturmuştur. CRISPR hedef dışı tahmininde sgRNA--hedef ilişkisini çizge bağlantısı olarak ele alan önceki çalışmalar, link prediction temelli GNN formülasyonunun bu problem alanına uygulanabilir olduğunu göstermiştir (Vinodkumar ve diğ., 2021). Bu tezde ise sgRNA--target bağlantı çizgesinin ötesine geçilerek, hedef gözlemine ait epigenetik ve kromatin bağlamının model içinde hangi temsille daha anlamlı kullanılabildiği sınanmıştır.

Çalışmanın değerlendirme ekseni ikiye ayrılmıştır. İlk eksen, eşikten bağımsız sıralama başarımıdır. Bu eksende area under the precision--recall curve (AUPRC) birincil metrik olarak kullanılmıştır. Dengesiz sınıf dağılımlarında precision--recall eğrileri, ROC eğrilerine göre pozitif sınıf başarımını daha doğrudan yansıtır (Davis ve Goadrich, 2006; Saito ve Rehmsmeier, 2015). İkinci eksen, validation üzerinde seçilmiş karar eşiğinde nadir ölçülmüş negatiflerin ayırt edilebilmesidir. Bu eksende Matthews correlation coefficient (MCC), macro-F1, specificity ve karışıklık matrisi değerleri ana yorum kanıtları olarak kullanılmıştır. Bu iki eksen aynı sonucu ölçmez. AUPRC modelin genel sıralama kalitesini değerlendirirken, karar eşiğine bağlı metrikler belirli bir kullanım noktasındaki sınıf ayırma davranışını gösterir.

Bu ayrım, çalışmanın ana iddia sınırını belirler. Bu tezde geliştirilen context-aware Graph C / GATv2 çizgisi, XGBoost F4 tabular referansına karşı sağlam bir AUPRC üstünlüğü olarak sunulmamaktadır. Bu iddia sınırı, Bölüm~\ref{chap:deneysel}'de verilen sağlamlık sonuçları ve paired-difference dağılımlarıyla desteklenmiştir (Çizelge~\ref{tab:auprc-saglamlik}; Şekil~\ref{fig:saglamlik-auprc-araliklari}; Şekil~\ref{fig:paired-difference-distributions}). Buna karşılık Graph C target-observation temsili, family-aware target-context encoder ve sequence/context fusion varyantları, validation-kilitli karar eşiğinde nadir ölçülmüş negatiflerin tanınması tarafında savunulabilir bir katkı göstermiştir. Bu katkının deneysel karşılığı Graph C GATv2, mekanizma çıkarma ve family-aware/context etkileşim sonuçlarında verilmiştir (Çizelge~\ref{tab:attention-sonuclari}; Çizelge~\ref{tab:graphc-mekanizma}; Çizelge~\ref{tab:family-aware-context}). Bu nedenle çalışma, en güçlü genel tahmin edici iddiası yerine, bağlam duyarlı çizge temsilinin nerede işe yaradığını ve nerede sınırlı kaldığını gösteren kontrollü bir modelleme ve değerlendirme çalışması olarak konumlandırılmıştır.

\section{Tezin Amacı}

Bu tezin amacı, CRISPR-Cas9 hedef dışı tahmininde epigenetik ve kromatin bağlamının çizge tabanlı modeller içinde temsil edilme biçimini incelemek ve bu temsilin hem sıralama başarımı hem de nadir negatif sınıf tanıma davranışı üzerindeki etkisini değerlendirmektir. Çalışmada Mak ve diğ. (2022) veri çizgisinden gelen crisprSQL türevi epigenetik-nükleozom veri seti kullanılmıştır. Ana hedef, aynı değerlendirme sözleşmesi altında tabular baseline, sequence tabanlı baseline, GCN, GAT/GATv2 ve Graph C target-observation temsillerinin karşılaştırılmasıdır.

Değerlendirme sözleşmesi, sonuçların güvenilir yorumlanabilmesi için dondurulmuştur. Ana ikili etiket Scheme A olarak tanımlanmış ve \texttt{cleavage\_freq > 1e-5} kuralı kullanılmıştır. NaN \texttt{cleavage\_freq} satırları supervised label üretiminden çıkarılmıştır. Ana veri ayrımı guide-disjoint olacak biçimde kurulmuştur. Ana validation ve test evreni measured-only satırlardan oluşmuş, \texttt{experiment\_id = 18} ana değerlendirme dışında bırakılmıştır. Ön işleme istatistikleri eğitim verisinden öğrenilmiş, model seçimi validation AUPRC ile yapılmış ve test seti üzerinde threshold veya model ayarı yapılmamıştır.

Bu çalışma kapsamında üç temsil düzeyi karşılaştırılmıştır. Graph A, fiziksel hedef düğümlerine dayanan minimal sgRNA--target çizgesi olarak kullanılmıştır. Graph B, rehber benzerliğiyle zenginleştirilmiş sınırlı bir kontrol olarak ele alınmıştır. Graph C ise her satırı ayrı bir target-observation düğümü olarak temsil ederek, aynı fiziksel hedef koordinatının farklı deneysel veya epigenetik bağlamlarda farklı gözlemler taşıyabilmesini model içine yansıtmıştır. Bu ayrım, Graph C sonucunun salt topology enrichment olarak yorumlanmaması için özellikle korunmuştur.

Tezin pratik amacı, güçlü bir tabular referans olan XGBoost F4 karşısında GNN modellerin koşulsuz üstün olduğunu göstermek değildir. Daha dar ve savunulabilir amaç, bağlam duyarlı Graph C GATv2 çizgisinin hangi koşullarda yararlı sinyal ürettiğini, hangi metriklerde sınırlı kaldığını ve bu etkinin hangi feature ailelerine bağlı olduğunu göstermektir. Bu nedenle AUPRC birincil metrik olarak korunmuş, MCC, macro-F1, specificity ve TN/FP/FN/TP değerleri ise nadir negatif operating-point davranışının ana kanıtları olarak değerlendirilmiştir.

\section{Literatür Araştırması}

CRISPR-Cas9 hedef dışı tahmin literatüründe ilk ana eksen, rehber RNA ve hedef DNA sekansları arasındaki benzerliğe dayalı modellemedir. Hsu ve diğ. (2013), Cas9 hedefleme özgüllüğünün rehber--hedef uyumsuzluk örüntülerinden güçlü biçimde etkilendiğini göstermiştir. Daha sonraki derin öğrenme yaklaşımları, sekans çiftlerini CNN, RNN, attention, çoklu görünüm veya transformer tabanlı temsilcilerle işleyerek hedef dışı aktiviteyi tahmin etmeye çalışmıştır (Chuai ve diğ., 2018; Lin ve diğ., 2020; Liu ve diğ., 2020; Luo ve diğ., 2024; Sun ve diğ., 2024). Bu çalışmalar, sekans temsilinin problem için zorunlu bir temel olduğunu göstermektedir. Ancak farklı veri kaynakları, negatif örnekleme stratejileri, guide split politikaları ve sınıf prevalansları nedeniyle bu literatürdeki ham skorlar bu tezdeki measured-only test evreniyle doğrudan karşılaştırılamaz.

İkinci eksen, epigenetik ve kromatin bağlamının hedef dışı aktiviteyle ilişkisidir. crisprSQL, farklı CRISPR/Cas hedef dışı kesim deneylerini yapılandırılmış biçimde bir araya getirmiştir (Störtz ve Minary, 2021). Mak ve diğ. (2022), bu veri çizgisi üzerinde 6 deneysel epigenetik özellik ve 13 hesaplanmış nükleozom ilişkili özelliği incelemiştir. piCRISPR çalışması da aynı veri ailesine yakın bir hatta fiziksel olarak bilgilendirilmiş sequence, kromatin ve nükleozom tanımlayıcılarının hedef dışı kesim tahminindeki rolünü incelemiştir (Störtz ve diğ., 2023). Bu tez, aynı veri ve feature lineage'ından yararlansa da Mak ve diğ. (2022) çalışmasının birebir reprodüksiyonu değildir. Mak çalışması continuous CA dönüşümü, korelasyon ve SHAP temelli feature analizi etrafında yapılandırılmıştır. Bu tezde ise Scheme A ikili sınıflandırma hedefi, guide-disjoint split, measured-only validation/test evreni ve AUPRC birincil metriği kullanılmıştır.

Üçüncü eksen, CRISPR hedef dışı tahmininin çizge problemi olarak ele alınmasıdır. GCN mimarisi, komşuluk ilişkilerinden düğüm temsili öğrenmek için temel bir GNN yaklaşımı sağlamıştır (Kipf ve Welling, 2017). Vinodkumar ve diğ. (2021), sgRNA hedef dışı aktivitesinin GCN tabanlı bir yaklaşımla modellenebileceğini göstermiştir. Bu tezde bu çizgi, daha katı bir measured-only ve guide-disjoint değerlendirme sözleşmesiyle genişletilmiştir. Minimal fiziksel hedef çizgesi, guide-similarity kontrolü ve target-observation context çizgesi ayrı ayrı test edilmiştir. Böylece graph schema seçiminin ve hedef bağlamı semantiğinin model davranışı üzerindeki etkisi ayrıştırılmıştır.

Dördüncü eksen, attention tabanlı GNN mimarileridir. GAT modelleri, komşu düğümlerden gelen mesajları attention katsayılarıyla ağırlıklandırır (Veličković ve diğ., 2018). GATv2, attention mekanizmasının daha dinamik bir biçimini hedeflediği için Graph C bağlamında anlamlı bir model adayıdır (Brody ve diğ., 2022). Ancak bu tezde attention mekanizması tek başına başarı garantisi olarak ele alınmamıştır. Attention tabanlı modeller Graph A/B/C şemaları ve edge-aware kullanım koşulları altında kontrollü biçimde sınanmıştır. Graph C target-observation bağlamı ve candidate-edge enerji özelliklerinin birlikte kullanılması, nadir negatif operating-point davranışını güçlendiren ana yön olarak değerlendirilmiştir.

Beşinci eksen, değerlendirme metodolojisi ve sınıf dengesizliği yorumudur. Precision--recall eğrileri, sınıf dağılımının dengesiz olduğu durumlarda model karşılaştırmasını ROC eğrilerine göre daha anlamlı hale getirebilir (Davis ve Goadrich, 2006; Saito ve Rehmsmeier, 2015). Precision değerinin sınıf prevalansına duyarlı olması, AUPRC değerlerinin test evrenindeki sınıf oranıyla birlikte yorumlanmasını gerektirir (Williams, 2021). Bu tezde test evreni pozitif ağırlıklıdır ve pozitif prevalans 0.900705'tir. Bu nedenle problem, genom çapında negatif ağırlıklı retrieval benchmark'larından farklıdır. Nadir sınıf bu ölçüm evreninde negatiftir. Bu durum, specificity, MCC, macro-F1 ve karışıklık matrisinin yorum değerini artırmıştır. Sağlamlık analizinde guide-cluster bootstrap, paired-difference analizleri ve çoklu tohum değerlendirmesi kullanılarak tekil eğitim koşulundaki AUPRC kazanımlarının belirsizlik içinde kalıp kalmadığı ayrıca incelenmiştir (Field ve Welsh, 2007; Bethard, 2022).

Bu literatür içinde tezin konumu, sekans odaklı bir karşılaştırma çalışması veya Mak ve diğ. (2022) reprodüksiyonu değildir. Çalışma, epigenetik/kromatin bağlamının çizge temsiliyle nasıl kullanılabileceğine odaklanır. Bu nedenle literatürden alınan fikirler birebir yeniden üretim iddiasıyla değil, yerel veri sözleşmesi altında uyarlanmış modelleme kararları olarak ele alınmıştır.

\section{Hipotez}

Bu tezde üç temel hipotez sınanmıştır.

Birinci hipotez, Graph C target-observation context temsili ve family-aware GATv2 encoder'ın, sabit measured-only ve guide-disjoint değerlendirme sözleşmesi altında nadir negatif sınıf tanıma davranışını Graph A GCN ve düz Graph C GCN referanslarına göre iyileştireceğidir. Bu hipotez, AUPRC yerine validation-kilitli threshold metrikleriyle ilişkilidir. MCC, macro-F1, specificity ve TN/FP/FN/TP değerleri bu davranışın ana kanıtları olarak değerlendirilmiştir. Hipotezin deneysel karşılığı Graph C GATv2 ve family-aware encoder sonuçlarında izlenmiştir (Çizelge~\ref{tab:attention-sonuclari}; Çizelge~\ref{tab:family-aware-context}; Şekil~\ref{fig:sequence-context-threshold}).

İkinci hipotez, context-aware GNN çizgisinin XGBoost F4'e karşı sağlam bir AUPRC üstünlüğü oluşturmayacağıdır. Bu hipotez, çalışmanın iddia sınırını açık biçimde belirler. XGBoost F4, ilk dondurulmuş değerlendirmede test AUPRC 0.992522 ile en güçlü tabular referans olarak kalmıştır (Çizelge~\ref{tab:non-graph-referanslar}). Sağlamlık analizinde XGBoost version drift nedeniyle yeniden üretilen F4 AUPRC değeri 0.992338 olarak raporlanmış, paired AUPRC farkları ise sıfırı dışlamamıştır (Çizelge~\ref{tab:auprc-saglamlik}; Şekil~\ref{fig:paired-difference-distributions}). Bu nedenle GNN sonuçları rekabetçi fakat AUPRC üstünlüğü iddiası taşımayan sonuçlar olarak yorumlanmıştır.

Üçüncü hipotez, Graph C kazanımının explicit context-similarity edge'lerinden çok doğrudan target-observation context feature'larına, özellikle experimental epigenetic feature ailesine bağlı olduğudur. Mekanizma ve alt-özellik maskeleme analizleri bu hipotezin temel deneysel karşılığıdır. Bu kanıtlar Çizelge~\ref{tab:graphc-mekanizma}, Çizelge~\ref{tab:target-context-mask} ve Şekil~\ref{fig:target-context-subgroup-threshold}'de verilmiştir. Target-observation context feature'ları maskelendiğinde model davranışının belirgin biçimde bozulması, Graph C katkısının ek topolojiyle tek başına açıklanamayacağını göstermiştir. Bu bulgu biyolojik olarak kromatin ve epigenetik bağlam literatürüyle uyumludur; ancak attention ağırlıkları, feature masking sonuçları veya context gate davranışları nedensel biyolojik kanıt olarak yorumlanmamıştır.

Bu hipotezler birlikte ele alındığında tez, CRISPR-Cas9 hedef dışı tahmininde context-aware GNN yaklaşımının değerini birincil AUPRC üstünlüğüyle değil, kontrollü değerlendirme altında mekanizma izole edilmiş ve validation-kilitli nadir negatif tanıma davranışıyla ortaya koymayı amaçlamaktadır.
